% ============================================================================
% Fichier  : annexes/Ann_D_textes.tex
% Titre    : Textes et Normes de Référence
% Version  : 1.0 -- novembre 2024
% ============================================================================
% Les extraits de la loi n° 2010/012 sont recopiés du texte officiel détenu
% dans COURS\Normes\loi_2010-012_cybersecurite_cybercriminalite.pdf (lu le
% 08/11/2024).
% ============================================================================

\chapter{Textes et Normes de Référence}
\label{ann:textes}

Cette annexe donne, pour chaque texte cité, son degré de vérification.
Trois degrés sont utilisés~: \emph{texte officiel lu} (article cité entre
guillemets d'après le texte officiel)~; \emph{paraphrasé} (norme sous
licence, dont le contenu est résumé sans citation)~; \emph{à vérifier}
(information dont l'auteur n'a pas pu contrôler la source à la date de
rédaction).

\section*{Droit camerounais}

\begin{tableaucours}{L{3.8cm}L{5.2cm}L{3.2cm}}{%
    \tableauHeader{Texte} &
    \tableauHeader{Articles utilisés} &
    \tableauHeader{Degré}
}
Loi n°~2010/012 du 21~décembre 2010 relative à la cybersécurité et à la
cybercriminalité & 26 al.~4, 29 al.~1 à~3, 30 al.~1, 52, 53, 54, 55, 56,
65 al.~2, 69, 72, 91 al.~3 & Texte officiel lu \\
Loi n°~2010/012 & 26 al.~3, 57, 58, 66, 68, 71 & Texte officiel lu, cités
par résumé \\
Loi n°~2024/017 du 23~décembre 2024 (données personnelles) & Mentionnée
seulement & Renvoi à l'UE JUR01 \\
Code de procédure pénale & Preuve, expertise & À vérifier \\
Réglementation de la supervision des établissements de microfinance &
Notification d'incident & À vérifier \\
\end{tableaucours}

\section*{Extraits de la loi n°~2010/012 non reproduits dans les chapitres}

\begin{boxjuris}
\textbf{Article~66, alinéa~1}~: \textit{« Est puni d'un emprisonnement de
deux (02) à cinq (05) ans et d'une amende de 1.000.000 (un million) à
2.000.000 (deux millions) F CFA ou de l'une de ces deux peines seulement,
celui qui entraîne la perturbation ou l'interruption du fonctionnement d'un
réseau de communications électroniques ou d'un équipement terminal, en
introduisant, transmettant, endommageant, effaçant, détériorant, modifiant,
supprimant ou rendant inaccessibles les données. »}

\medskip
\textbf{Article~68, alinéa~1}~: \textit{« Est puni d'un emprisonnement de
cinq (05) à dix (10) ans et d'une amende de 10.000.000 (dix millions) à
50.000.000 (cinquante millions) F CFA ou de l'une de ces deux peines
seulement, celui qui accède ou se maintient, frauduleusement, dans tout ou
partie d'un réseau de communications électroniques ou d'un système
d'information en transmettant, endommageant, provoquant une perturbation
grave ou une interruption du fonctionnement dudit système ou dudit
réseau. »}

\medskip
\textbf{Article~71}~: \textit{« Est puni d'un emprisonnement de deux (02) à
cinq (05) ans et d'une amende de 1.000.000 (un million) à 25.000.000 (vingt
cinq millions) F CFA, celui qui introduit sans droit, des données dans un
système d'information ou dans un réseau de communications électroniques en
vue de supprimer ou de modifier les données qui en sont contenues. »}
\end{boxjuris}

\section*{Conventions internationales}

\begin{tableaucours}{L{4.6cm}L{4.2cm}L{3.4cm}}{%
    \tableauHeader{Texte} &
    \tableauHeader{Objet} &
    \tableauHeader{Degré}
}
Convention de Budapest (Conseil de l'Europe, 2001) & Cybercriminalité,
conservation rapide, entraide & Statut du Cameroun à vérifier \\
Convention de Malabo (Union africaine, 2014) & Cybersécurité et données
personnelles~; en vigueur en 2023 & Statut du Cameroun à vérifier \\
\end{tableaucours}

\section*{Normes et référentiels}

\begin{tableaucours}{L{3.6cm}L{5.6cm}L{3.0cm}}{%
    \tableauHeader{Texte} &
    \tableauHeader{Objet} &
    \tableauHeader{Degré}
}
ISO/IEC 27037:2012 & Identification, recueil, acquisition, préservation &
Paraphrasé \\
ISO/IEC 27041:2015 & Adéquation des méthodes & Paraphrasé \\
ISO/IEC 27042:2015 & Analyse et interprétation & Paraphrasé \\
ISO/IEC 27043:2015 & Principes et processus d'investigation & Paraphrasé \\
ISO/IEC 27035 & Gestion des incidents & Paraphrasé \\
ISO/IEC 27002:2022, mesure 8.8 & Gestion des vulnérabilités techniques &
Paraphrasé \\
ISO/IEC 17025:2017 & Compétence des laboratoires & Paraphrasé \\
ISO/IEC 29147, 30111 & Divulgation et traitement des vulnérabilités &
Paraphrasé \\
NIST SP 800-86, 800-61 r2, 800-115, 800-40 r4 & Forensique, incidents,
évaluation, correctifs & Publications libres, paraphrasées \\
RFC 3227 & Ordre de volatilité & Publication libre \\
CVSS v4.0 et v3.1 (FIRST) & Notation de la gravité & Spécification libre \\
EPSS (FIRST), KEV et SSVC (CISA) & Probabilité, exploitation, décision &
Ressources libres \\
\end{tableaucours}
